% auto-generated by extract_demographics.py -- do not hand-edit
\textbf{Field} & \textbf{LIDC} & \textbf{Aggregate} \\
\midrule
N (patients, unique) & 1010 & 1010 \\
Age, median (IQR) [yr] & 61 (50--70) & 61 (50--70) \\
Age, range [yr] & 14--88 & 14--88 \\
Sex (M / F), n (\%) & 140 / 140 (28\%) & 140 / 140 (28\%) \\
BMI, median (IQR) & N/A & N/A \\
Scan-year range & N/A & N/A \\
% --- caption (paste into manuscript.tex, replacing the old one) ---
\caption{Patient-level demographics. v0.5 contains LIDC only (n = 1{,}010); the AAPM/Mayo columns are the v1.0 roadmap plan and are not part of v0.5. The LIDC cells are regenerated from \texttt{metadata.json} by the committed extraction script \texttt{analysis/extract\_demographics.py} (build provenance: 1010 files, combined SHA-256 \texttt{20476a8105319e6192eebf233e274a4f6f5f787a574bd605809963f3160df870}). $^{*}$LIDC age and sex are recorded only on a subset of patients (age 20\%, sex 28\%) per LIDC's variable DICOM-header preservation; the percentages in the Sex row are computed over the patients with known sex. $^{\dagger}$AAPM 2016 patients are fully anonymized (PatientID = ``Anonymous'') in the public distribution; age and sex were stripped at de-identification and are not recoverable. BMI and StudyDate are not preserved in any of the source releases.}
